% Generated by `taskgen course examples` from factoring.ec_mod_n_add (seed 5).
% Do not edit: change examples.yaml and regenerate.
\begin{example}[сложение точек над $\mathbb{Z}/N\mathbb{Z}$]\label{ex:ecm-add-mod-n}
Кривая $y^2 = x^3 + 396x + 444$ рассматривается над $\mathbb{Z}/493\mathbb{Z}$. Вычислить $P + Q$ для $P = (103, 318)$, $Q = (174, 322)$ по обычным формулам с делением по модулю $493$. Если деление невозможно, найти нетривиальный делитель числа $493$.

\par\noindent\textit{Решение.} Точки лежат на кривой (проверка подстановкой). Сложение: $\lambda = \dfrac{y_2 - y_1}{x_2 - x_1}$, знаменатель $d = 174 - 103 \equiv 71$, числитель $322 - 318 \equiv 4$.  Обратный к $71$ по модулю $493$ (расширенный алгоритм Евклида для $(493, 71)$):
\[\begin{array}{r|r|r|r}
r_i & q_i & u_i & v_i \\ \hline
493 &  & 1 & 0 \\
71 &  & 0 & 1 \\
67 & 6 & 1 & -6 \\
4 & 1 & -1 & 7 \\
3 & 16 & 17 & -118 \\
1 & 1 & -18 & 125 \\
0 & 3 & 71 & -493 \\
\end{array}\qquad r_i = u_i \cdot 493 + v_i \cdot 71\]

$71^{-1} \equiv 125\ (\textrm{mod}\ 493)$.
$\lambda = 4 \cdot 125 \equiv 7$, $x_3 = \lambda^2 - x_1 - x_2 \equiv 265$, $y_3 = \lambda(x_1 - x_3) - y_1 \equiv 27$.\[P + Q = (265, 27).\]\par Проверка по модулю простых делителей ($493 = 17 \cdot 29$): приведения ответа $(10, 10) \in E(\mathbb{F}_{17})$ и $(4, 27) \in E(\mathbb{F}_{29})$ --- это суммы приведений $P$ и $Q$ (приведение по модулю простого делителя сохраняет операцию).
\end{example}
